% Expository adaptation, October 7, 2026. Original mathematics: OpenAI, September 23, 2026.
% Prose revised and explanations expanded; see README.md for the source and scope.
\section{Sparse random tables for the exact construction}
\label{sec:atom}

The hierarchy defines large matrices that we cannot store. The next step is to represent a requested entry by the conditional expectation of a sparse random table. In one environment, a row of such a table has only a bounded number of nonzero entries. Averaging over environments recovers the exact contribution we need.

There are three kinds of contributions: transition and residual entries, products involving corrections, and column detector scores. The tables called bins below separate entries by magnitude; correction tables permit local movement in both directions; detector tables supply the scores used to decide which columns permit correction. A normalized table also needs an earlier preparatory table to compute the probability by which it divides. These are roles within one construction, not separate approximation algorithms.

We first prove the exact conditional-mean identities.  We then extend the
same row formulas to computed raw arrays and prove the support and access
bounds that hold even when those arrays are inaccurate.  Finally we bound
the effect of errors in the incoming gates and column multipliers.  The
later implementation supplies enumeration and pairing of the resulting
ports.

Exact conditional-mean identities refer only to the true tables defined
here.  Applying the same formulas to approximate raw signals gives
approximate tables; their normalization need not be unbiased.

Weighted pseudorandomness provides a related way to represent expectations
with few random bits; see Braverman, Cohen and Garg~\cite{BravermanCohenGarg2020}
and the error-reduction work of Cohen, Doron, Renard, Sberlo and
Ta-Shma~\cite{CohenDoronRenardSberloTaShma2021}. The tables below have
exact conditional-mean identities. Their estimates use a common
environment, so products and subsequent normalization are analyzed
through joint error bounds rather than independent sampling.

\subsection{Conditioning and normalization}

Throughout this section $s=8$.  Fix a stage, write $f=b_l$, and let $G$ be
one of its nonnegative raw matrix inputs.  The hierarchy gives absolute
constants $M_G,K_G$ such that
\begin{equation}
 \sum_y G(x,y)\le M_G,\qquad G(x,y)\le K_G f.
 \label{eq:atom-input-bounds}
\end{equation}
These bounds include pilot and actual inputs and the residual $D$; see
Lemmas~\ref{lem:hier-pilot-bounds} and~\ref{lem:hier-residual}.
For a correction $F=fQ(G/f)g_0$, the additional input is
\begin{equation}
 g_0(y)>0\quad\Longrightarrow\quad
 \#\{x:G(x,y)\ge\delta_1 f\}\le R_*/f,
 \label{eq:atom-column-input}
\end{equation}
from Lemma~\ref{lem:hier-column-test}.

A row task has start rate $u\le f$, end scale $a$, and, when needed, a
usage rate $t\le a$, all on the rate grid.  Its domain is
$X_u(x)=Y_a(y)=1$, with distinct start and terminal channels.  We use
$a\le K_f f$ for an absolute $K_f$, and use $a=f$ for corrections and
detector scores.  Here and in the estimation and compression sections,
$t$ denotes a sampling rate; the numerical clock $t$ is introduced
separately in Section~\ref{sec:arith}.  The true first raw signal is
\begin{equation}
 d_y=G(x,y)/a.
 \label{eq:atom-first-signal}
\end{equation}
An optional second raw signal is
$\zeta\mathbf1_{\{Y_t(y)=1\}}q_y$, where $q_y=q(x,y)\in[0,1]$ is a
deterministic conditional mean and $\zeta>0$ is a fixed sufficiently small
constant.  That mean will be specified separately for each normalized task.

A query must refer to the same table regardless of the path by which the queried entry was found. Otherwise repeated contributions to one entry could not be combined consistently. For this reason, all parameters of a table are common to the whole table: its stage and
iteration, task type, channel roles, channel padding suffixes, rates
$u,a,t$, and subsequently its budget and numerical convention.  They do
not depend on the row or column at which the table is queried.

Write $\operatorname{rk}_k$ for the $k$th largest value, with unlimited
zero padding.  For two arrays on the same padded index set,
\begin{equation}
 |\operatorname{rk}_k(v)-\operatorname{rk}_k(w)|
 \le\sup_z|v_z-w_z|.
 \label{eq:atom-rank-lipschitz}
\end{equation}
Indeed, if every coordinate changes by at most $\epsilon$, at least $k$
coordinates at or above the old $k$th value remain at or above that value
minus $\epsilon$; apply this in both directions.

Here is the basic tradeoff. A gate can force every sampled row to be sparse by setting crowded rows to zero. That deletion changes the expectation. Dividing by the gate's true conditional mean corrects the expectation, provided the mean is positive and not too small. The normalization can be seen before choosing the row functions.  Fix a
pair $(x,y)$ on its conditioning slice, a deterministic nonnegative value
$k(d_y)$, and a random row gate $R\in[0,1]$.  Choose a preparatory factor
$h$ that is one wherever $k$ is positive.  If
\[
 q(x,y)=\mathbb E[h(d_y)R\mid X_u(x)=Y_t(y)=1]\ge\tfrac12
\]
on that support, then $k(d_y)R/q(x,y)$ has conditional mean $k(d_y)$
there.  A clipped denominator keeps the formula bounded at every other
position and for computed inputs.  The constructions below arrange a
uniform positive lower bound on this mean while the row gate limits the
number of nonzero retained entries.  Corrections use an additional incoming
gate, with the same normalization principle.

We will use several continuous switches to distinguish small from large entries. The numerical values of their thresholds are less important than the gaps between them: an entry can be safely included in an access list before it becomes nonzero, and small numerical perturbations cannot destroy that inclusion. The next definitions fix one concrete set of switches and margins so that these requirements can be checked exactly.

For real $\alpha<\beta$, let $\chi_{\alpha,\beta}$ be a fixed increasing
switch, zero on $(-\infty,\alpha]$ and one on $[\beta,\infty)$.
Choose it by rescaling the normalized integral of
$z^{s-1}(1-z)^{s-1}$ on $[0,1]$.
Both $\chi_{\alpha,\beta}^{1/s}$ and
$(1-\chi_{\alpha,\beta})^{1/s}$ are Lipschitz: at either endpoint the
vanishing function has a zero of order $s$, and on the interior ordinary
differentiation applies.  Products of bounded functions with Lipschitz
$s$th roots again have Lipschitz $s$th roots.

Every scalar saturation below uses fixed min/max clipping,
\[
 \operatorname{clip}_{[\alpha,\beta]}(v)
   =\min\{\beta,\max\{\alpha,v\}\},
\]
which is $1$-Lipschitz.  We specify its bounds when defining the function.
Functions initially defined for nonnegative first entries are extended
by zero below zero; each already vanishes below a positive support
threshold.  These conventions apply only to evaluation of the row
functions.  Ranks and access tests always use the full prescribed raw
arrays before scalar clipping.

Here is one explicit set of margins used below.  Let $\vartheta>0$ be a
fixed lower support threshold of a main row function, and put
$\rho=\operatorname{rk}_K(v)$ for its first raw row $v$.  Define
\begin{align}
 h_{\vartheta}(v)
   &=\chi_{\vartheta/2,\,3\vartheta/4}(v),\notag\\
 R_{\vartheta}(\rho)
   &=1-\chi_{\vartheta/1000,\,\vartheta/500}(\rho),\notag\\
 I_{\vartheta}(v,\rho)
   &=\chi_{\vartheta/8,\,\vartheta/4}(v)
       \bigl(1-\chi_{\vartheta/100,\,\vartheta/50}(\rho)\bigr).
 \label{eq:atom-threshold-template}
\end{align}
We take $\vartheta=\theta_1$ for corrections, $\vartheta=1/2$ for dense
bins, and $\vartheta=\delta_0$ for the detector score.  For that detector
we make the allowed hierarchy choices
\[
 \eta_-=\delta_0/1000,\qquad \eta_+=\delta_0/500,\qquad
 \rho(t)=1-\chi_{\delta_0/1000,\,\delta_0/500}(t)
          =R_{\delta_0}(t).
\]
Here $\rho$ denotes the cutoff in the hierarchy, and the same switch
$\chi$ is used in both occurrences.  These choices satisfy
$0<\eta_-<\eta_+<\delta_0/3$.  Thus the detector row gate here is exactly
$R_i$ in Equation~\eqref{eq:hier-pilot-row-gate}.
The detector does not require $h_{\vartheta}$.
The hierarchy's choices may and do satisfy
$\delta_1<\theta_1/100$ and $0<\delta_0<\delta_1$.
Thus the correction auxiliary score $I_{\theta_1}$ can be positive only
when $G(x,y)>\delta_1 f$.

\subsection{Separate entries by size and normalize each part}

A single sampling rate is not well suited to entries of very different sizes. We split an entry into overlapping contributions at geometric scales. The smooth weights sum to one, so this introduces no approximation. At each scale, sampling and a row gate leave bounded row support; the preparatory conditional mean corrects for the gate's deletion.

Put $q_*=2^H$, and let $\chi=\chi_{1/2,1}$.  For $w\ge0$ define
\begin{align}
 \omega_0(w)&=\chi(w),\notag\\
 \omega_k(w)&=\chi(w/b_k)-\chi(w/b_{k-1})\qquad(k\ge1).
 \label{eq:atom-bin-weights}
\end{align}
The top bin has no preceding switch.  The weights are nonnegative and
sum to one for $w>0$, since their partial sums telescope to
$\chi(w/b_k)\to1$.  Hence
\begin{equation}
 w=\sum_{k\ge0}w\omega_k(w).
 \label{eq:atom-bin-partition}
\end{equation}
At scale $a=b_k$, the bin in atom units is
\begin{equation}
 k_a(d)=d\omega_k(ad).
 \label{eq:atom-bin-function}
\end{equation}
For $k\ge1$ its support is contained in $1/2<d<q_*$.
Its two switch transition intervals are disjoint, and
\[
 \chi(d)-\chi(d/q_*)=\chi(d)(1-\chi(d/q_*)),
\]
which proves directly that its $s$th root is Lipschitz.
For $k=0$, fix $B_{\rm top}>\max\{K_G,1\}$ and evaluate the defining
expression at $\operatorname{clip}_{[0,B_{\rm top}]}(d)$.
This makes the top-bin function bounded with Lipschitz root.  It
preserves every true value, since $a=b_0=1$ and
$G(x,y)\le K_G f\le K_G$.  For $k\ge1$, retain the zero extension
outside the support $(1/2,q_*)$; no upper clipping changes a zero bin
value.  Bins with $a>2K_G f$ vanish on the true input, so a fixed
$K_f\ge2K_G$ suffices for the permitted scale range.

For the routed residual bin use instead
\begin{equation}
 k_a^{(j)}(d)=k_a(d)s_j(ad/f).
 \label{eq:atom-routed-bin}
\end{equation}
The routing fractions have bounded Lipschitz roots.  Since $a/f\le K_f$
and the bin is zero above a fixed argument range (or is top-bin clipped),
these functions have uniformly bounded Lipschitz roots as well.  The
same zero-extension and clipping conventions apply to routed bins.
Summing Equation~\eqref{eq:atom-routed-bin} over scales gives the routed
residual contribution exactly on every true input.

Having split the magnitude scales, we now apply the normalization principle to one bin. Choose $R_{1/2}$ and $h=h_{1/2}$ from
Equation~\eqref{eq:atom-threshold-template}.  For the true row set
\[
 \rho_x=\operatorname{rk}_K
   \{G(x,z)/a:Y_a(z)=1\},\qquad R=R_{1/2}(\rho_x).
\]
Write
\begin{equation}
 q_y=\mathbb E\bigl[h(d_y)R\mid X_u(x)=Y_t(y)=1\bigr],
 \qquad
 c(q)=\min\{1,\max\{1/2,q\}\}.
 \label{eq:atom-dense-denominator}
\end{equation}
The number $q_y$ is the deterministic conditional mean at the indicated
pair.  The second raw signal carries the $Y_t$ membership mask.  In
addition, $q_y=0$ if $d_y\le1/4$.
The true preparatory table on the broad
terminal domain $Y_a$ is the unnormalized value $h(d_y)R$.  It uses only
the first raw signal.  Its conditional mean on the narrower slice is
exactly the $q_y$ in Equation~\eqref{eq:atom-dense-denominator}.
The usage rate $t$ is part of the fixed task tuple: changing the slice
on which this mean is taken would in general change the denominator.
The broad preparatory domain and the narrower normalization slice must
therefore be retained as separate parts of the interface.

\begin{lemma}[Dense atom identity]
\label{lem:atom-dense}
For sufficiently large fixed $K$, the table
\begin{equation}
 \mathcal B_a(x,y)=
  \mathbf1_{\{Y_t(y)=1\}}\frac{k_a(d_y)R}{c(q_y)}
 \label{eq:atom-dense-table}
\end{equation}
on retained starts has bounded values and row support.  On the slice
$X_u(x)=Y_t(y)=1$ its conditional expectation times $a$ equals
$G(x,y)\omega_k(G(x,y))$.  The same assertion holds for
$\mathcal B_a^{(j)}$ with the routing fraction included.
Apart from its start membership, its true value uses only the terminal
channel.
\end{lemma}

\begin{proof}
First show that the gate is usually fully open, even after conditioning on the queried destination. The row gate can fail to be full only if at least $K$ sampled columns
have $G(x,z)/a>r$, where $r=(1/2)/1000$ is fixed.
There are at most $M_G/(ar)$ such deterministic columns.
Conditional on $Y_t(y)=1$, the specified column contributes at most one
hit and each other column has $Y_a$-hit probability at most $C a$, by
Lemma~\ref{lem:env-ranks}.  Thus the expected count is at most
$1+CM_G/r$.  Conditioning also on the independent source channel does
not alter this bound.  Markov's inequality makes the failure probability
at most $1/8$ by a sufficiently large $K$.

Now the normalization is exact on every position that contributes to the bin. Where $k_a(d_y)>0$, we have $h(d_y)=1$.  Therefore
$q_y=\mathbb E[R\mid Y_t(y)=1]\in[7/8,1]$ there, and the clipping
does nothing.  The conditional mean of Equation~\eqref{eq:atom-dense-table}
is consequently $k_a(d_y)$.  Outside that support both sides vanish.
If a row output is positive, its first entry exceeds $1/2$ and its
$K$th first rank is below $1/1000$.  Hence there are at most $K-1$
positive entries.  Boundedness follows from the clipped denominator.  The routed
fraction is deterministic at truth and obeys the same argument.
\end{proof}

For products of a correction with a deterministic bin contribution, we
will use the following bounds.

\begin{lemma}[Bounds for a fine product contribution]
\label{lem:hier-bin-products}
Let $C$ be the transition matrix at the current stage, and let $F$ be
any finite pilot or actual correction there.  If a matrix $A_b\geq0$
on the same vertex set satisfies $A_b\mathbf 1\leq\mathbf 1$ and
$A_b(x,y)\leq c_b b$ for some $b>0$, then
\[
 (FA_b)\mathbf 1\leq k_*\mathbf 1,
 \qquad (FA_b)(x,y)\leq k_*c_b b.
\]
If also $A_b\leq C$, then $FA_b\leq FC$.  These estimates have constants
independent of the bin scale $b$.
\end{lemma}

\begin{proof}
Use $F\mathbf 1\leq k_*\mathbf 1$ from
Lemma~\ref{lem:hier-pilot-bounds}:
$(FA_b)\mathbf 1\leq F\mathbf 1$, and
$(FA_b)(x,y)\leq c_b b\sum_zF(x,z)$.  The last assertion follows by
multiplying $A_b\leq C$ on the left by the nonnegative matrix $F$.
\end{proof}

\subsection{Corrections with an incoming gate}

A correction is used as the left factor of a product. The algorithm moves along its edge, evaluates the right factor at the new vertex, and later returns. Both the implementation and its error analysis need a bounded number of incoming as well as outgoing accesses. A row gate alone provides no incoming bound.

We therefore add a column gate. The deterministic column-support bound proved for the hierarchy ensures that this additional gate is usually fully open on a true contributing entry. A union bound will show that both gates are open simultaneously with probability at least $3/4$; their independence is not required.

Consider $F=fQ(G/f)g_0(y)$ and put $a=f$, $\vartheta=\theta_1$.
Fix $B_Q>\max\{K_G,\theta_2\}$ and define the atom cutoff
\[
 Q_{\rm at}(d)=Q\bigl(\operatorname{clip}_{[0,B_Q]}(d)\bigr).
\]
It is bounded, has a Lipschitz $s$th root on the real line, and agrees
with $Q(d)$ for every true argument $d=G(x,y)/f$.  Its extension to
negative arguments is zero.  The exact hierarchy continues to use the
uncapped $Q$.  Let $R=R_{\theta_1}(\rho_x)$, $h=h_{\theta_1}$, and
$I=I_{\theta_1}$ be the functions in
Equation~\eqref{eq:atom-threshold-template}, where the true row rank is
$\rho_x=\operatorname{rk}_K\{G(x,z)/f:Y_f(z)=1\}$.
On a true or computed raw table, form $I(x,y)$ from that table's own row
entry and own rank.
With $\chi_{1/4,1/2}$ as above, define
\begin{equation}
 T_y=1-\chi_{1/4,1/2}\left(
       \operatorname{rk}_{K_c}\{I(x,y):X_u(x)=1\}\right).
 \label{eq:atom-incoming-gate}
\end{equation}
Thus $T_y=1$ through rank $1/4$ and $T_y=0$ from rank $1/2$.
This rank is taken over the fixed retained source set before any incoming
port filtering; there is no gate-dependent change of its index set.
The true correction preparation is $h(d_y)R T_y$ on $Y_a$.
Define its conditional mean by
\begin{equation}
 q_y=\mathbb E\bigl[h(d_y)R T_y\mid X_u(x)=Y_t(y)=1\bigr].
 \label{eq:atom-correction-denominator}
\end{equation}
The number $q_y$ is deterministic; as before, its second raw signal
carries the $Y_t$ membership mask.  Define the normalized true correction
using the clipped expression everywhere:
\begin{equation}
 \mathcal F(x,y)=\mathbf1_{\{Y_t(y)=1\}}
     \frac{Q_{\rm at}(d_y)R T_y}{c(q_y)}\,g_0(y).
 \label{eq:atom-correction-table}
\end{equation}
Here $q_y=0$ whenever $d_y\le\theta_1/2$.  Thus both kinds of second
signal are supported above a fixed positive first-signal threshold.
Choose $\zeta\le10^{-6}\min\{\theta_1,\delta_0,1/2\}$, decreasing it
later if a fixed compression margin requires this.  The second signals
are then uniformly tiny on the first-signal threshold scale.

\begin{lemma}[Correction atom identity, including zero multipliers]
\label{lem:atom-correction}
For sufficiently large fixed $K$ and then $K_c$, the conditional mean of
$f\mathcal F(x,y)$ on $X_u(x)=Y_t(y)=1$ is exactly $F(x,y)$.
On a true position with $Q(d_y)>0$ and $g_0(y)>0$,
\begin{equation}
 q_y=\alpha(x,y):=
  \mathbb E[R T_y\mid X_u(x)=Y_t(y)=1]\ge3/4.
 \label{eq:atom-alpha-bound}
\end{equation}
No lower bound on $q_y$ or $\alpha(x,y)$ is required when $g_0(y)=0$.
The true gates and normalization use only the two role channels $X,Y$;
$g_0$ is deterministic.
\end{lemma}

\begin{proof}
We bound failure of each gate separately. The row count argument in Lemma~\ref{lem:atom-dense}, with the correction
thresholds, gives $\Pr(R\ne1\mid\text{the two hits})\le1/8$.
If $g_0(y)>0$, every source with $I(x',y)>0$ belongs to the deterministic
set in Equation~\eqref{eq:atom-column-input}.  Conditional on $X_u(x)=1$,
the number of sampled members of that set has expectation at most
\[
 1+C uR_*/f\le1+C R_*.
\]
The bound remains valid after conditioning on the independent terminal
hit; it ignores the auxiliary row gate and therefore needs no additional
independence between rows.  If fewer than $K_c$ such sources are sampled,
the $K_c$th auxiliary rank is zero and $T_y=1$.  Choose
$K_c\ge8(1+C R_*)$ to make $\Pr(T_y\ne1\mid\text{the two hits})\le1/8$.
The union bound, without independence of the gates, gives
$\mathbb E[RT_y\mid\text{the two hits}]\ge3/4$.

On a contributing entry, divide by the mean of the product of the gates. Since $h(d_y)=1$ on $Q$'s support, Equation~\eqref{eq:atom-alpha-bound}
holds and $c(q_y)=q_y$ there whenever $g_0(y)>0$.  The claimed conditional
identity follows by division by this exact conditional mean.  If
$Q(d_y)=0$ or $g_0(y)=0$, both sides of the identity are zero.  In the latter
case the pre-incoming quantity $Q_{\rm at}(d_y)R/c(q_y)$ is still a well-defined
bounded comparison target, even if the unclipped mean vanishes.
\end{proof}

\subsection{Detector scores}

Detector scores already take values in a bounded interval, and their defining expectation does not need the preceding normalization. We only need to expose the row-bounded table whose column order statistic is the detector trial. The unnormalized detector table is
\begin{equation}
 \mathcal S_i(x,y)=\phi(G_i(x,y)/f)
 R_{\delta_0}\!\left(\operatorname{rk}_K
       \{G_i(x,w)/f:Y_f(w)=1\}\right)
 \quad\text{on }X_f(x)=Y_f(y)=1.
 \label{eq:atom-detector-table}
\end{equation}
Its column trial is the $J$th largest such score over retained starts.
By the identical cutoff choice, its conditional mean given $Y_f(y)=1$
is exactly $U_i(y)$ from Equation~\eqref{eq:hier-detector}.
The trial is an inverse rank access to a row-bounded table and does not
require an incoming degree bound.

\subsection{Accessing entries when the input estimates are inaccurate}

So far we have proved identities for true inputs. The algorithm must also have well-defined bounded access when its inputs are inaccurate, since the final deterministic enumeration visits inaccurate environments as well. The next bounds are therefore pointwise statements about arbitrary computed arrays, with no accuracy hypothesis.

The conditional-mean identities now give the true contributions.  A
product estimate will query a correction at $(x,z)$ and then a bin at
$(z,y)$.  Bounded outgoing degree limits the possible midpoints; bounded
incoming degree is needed when errors at the new source $z$ are summed
over old sources $x$.  We therefore specify access relations as well as
nonzero table supports.  The access relations may contain entries of
weight zero.

Passing through one copy projection is partitioned by its fixed source and
destination symbols.  Within each part the inherited hash argument
appends that symbol before the inherited suffix.  Lemma~\ref{lem:env-ranks}
then gives the same distribution as the unpadded channel.  The task tuple
records these fixed symbols and suffixes.  By definition, neither a right
factor's task parameters nor a column vector's task parameters retain an
earlier source address.
The environment supplied to the table is common to all its entries.
The vertices $x,y$ are arguments of that table, rather than additional
parameters.  Thus two paths that discover the same entry in that
environment query the same value; the path used to discover it is not
part of its mathematical definition.

We will apply these row functions to arrays represented by a finite list
of exceptional values and one common value at every omitted endpoint.
Call that common value the \emph{default}.  Adjoin arbitrarily many dummy
coordinates with the same default, and include them in every rank; they
are never matrix entries or ports.  For the true row arrays, dummy values
are zero.  Section~\ref{sec:buf-compression} constructs the finite-list
representations.  The elementary rank convention below lets us prove
support and access bounds before choosing that construction.

\begin{remark}[Ranks of arrays with a common default]
\label{rem:atom-finite-ranks}
Suppose a row array equals a common default outside a finite list,
and arbitrarily many dummy coordinates have that same default.  Its
$k$th rank can be computed from the listed coordinates and $k$ copies of
the default, with the stipulated zero padding.  Indeed, every omitted
coordinate merely supplies another copy of a value already available
with multiplicity at least $k$.  Additional listed real coordinates
equal to the default therefore do not change the rank.  The same
observation applies to a rank of entry sizes when each omitted entry
has the same default vector.

This is an equality of ranks for one prescribed full array.  The dummy
coordinates are never matrix entries or ports.  A common componentwise
clamp of the entries and their default preserves this representation:
an omitted position still equals the clamped default.  No assertion
that different calculations select identical lists or ports is needed.
\end{remark}

An access relation may safely contain an entry whose weight is zero. What matters is that it contains every nonzero entry, has bounded degree, and can be selected by repeatable finite-accuracy tests. The next lemma constructs such a relation by putting the access threshold below the nonzero threshold, with a gap large enough to absorb numerical error.

\begin{lemma}[Support margins and exposed row classes]
\label{lem:atom-margins}
Suppose a bounded nonnegative main function $k(v)$ vanishes for
$v\le\vartheta$.  Its product $k(v)R_{\vartheta}(\rho)$ and the
preparatory product $h_{\vartheta}(v)R_{\vartheta}(\rho)$ have at most
$K-1$ nonzero real entries in a row.  They admit a common exposed class
$A$, of at most $K-1$ real entries, containing their supports and satisfying
$I_{\vartheta}=1$ throughout $A$.  These assertions hold for arbitrary
real raw estimates, without a sign or accuracy hypothesis.

The auxiliary score $I_{\vartheta}$ also has at most $K-1$ nonzero real
entries and has an exposed row class with a bounded-support smooth
witness equal to one on that class.
\end{lemma}

\begin{proof}
The basic counting observation is that an entry threshold above the row's $K$th-largest value can be met by fewer than $K$ entries. For the main and preparatory products a nonzero entry has
$v>\vartheta/2$ and $\rho<\vartheta/500$.  There cannot be $K$ such
entries, since then the $K$th rank would exceed $\vartheta/2$.

We next leave sufficient room to perform the access test with finite precision. For an explicit class choose consistently computed approximations
$\widetilde v,\widetilde\rho$ with absolute errors at most
$\eta=\vartheta/10000$, and set
\begin{equation}
 A=\{y:Y_a(y)=1,\quad
       \widetilde v_y>3\vartheta/8,\quad
       \widetilde\rho<\vartheta/200\}.
 \label{eq:atom-exposed-class}
\end{equation}
Every positive main or preparatory entry belongs to $A$.  An entry in
$A$ satisfies $v_y>(3/8-1/10000)\vartheta>\vartheta/4$ and
$\rho<(1/200+1/10000)\vartheta<\vartheta/100$, so its witness
$I_{\vartheta}$ equals one.  Its entry threshold exceeds its rank
threshold, proving $|A|\le K-1$.

The support of $I_{\vartheta}$ requires $v>\vartheta/8$ and
$\rho<\vartheta/50$, again giving the same row bound.  To access this
score itself, the class
\[
 A_I=\{y:Y_a(y)=1,\quad
       \widetilde v_y>\vartheta/10,\quad
       \widetilde\rho<\vartheta/40\}
\]
contains its support and has at most $K-1$ entries.  A smooth unit witness
on $A_I$ is
\[
 J_{\vartheta}(v,\rho)=
  \chi_{82\vartheta/1000,\,90\vartheta/1000}(v)
  \bigl(1-\chi_{26\vartheta/1000,\,27\vartheta/1000}(\rho)\bigr).
\]
The approximation margins imply $A_I\subseteq\{J_{\vartheta}=1\}$.
Moreover $82/1000>3(27/1000)$, so this witness also has bounded row
support and the required entry versus rank separation.  The witness is
an analytical output; its support need not be broadened recursively.

These arguments use the rank of the prescribed entire raw array.
In particular, if at least $K$ dummy positions have a common virtual
default $\beta$, then $\rho\ge\beta$.  A real position still equal to
that default cannot meet either exposed class: the required lower bound
on $v=\beta$ exceeds the required upper bound on $\rho$.  Thus omitted
default positions do not create unlisted ports.
\end{proof}

The fixed-accuracy approximations in this lemma are part of the access
interface, not tests for equality of exact reals.  The implementation
must use the same approximation and tie conventions whenever it queries
the same table.  The strict margins prove correctness for every such
choice.  Auxiliary-only access classes require a row bound but no
incoming bound.
In particular, different computed versions may expose different
zero-weight entries.  The assertions are applied to each version's own
final raw array and repeatable approximation convention; they do not
require those classes to coincide.  Accuracy of the raw array is used
later in the error estimates, not in this support argument.


We can now specify the assembly order for an estimated table. First complete its raw row arrays. Apply the bounded row functions to those \emph{final} arrays, and only then form their incoming scores and gates. No unfinished intermediate row is used as though it were the final table.  In a normalized task its second input is a
computed value approximating $\zeta\mathbf1_{\{Y_t(y)=1\}}q_y$: the
calculator divides that input by $\zeta$ and applies $c$.  It does not launch another averaging
operation after selecting an endpoint.  Its auxiliary incoming scores
are formed from its own final first-signal rows, and its incoming gate
is Equation~\eqref{eq:atom-incoming-gate} for those scores.  A preparatory
task is an earlier, separate task; its raw arrays, scores, and gates are
its own, and its target is the unnormalized preparatory table specified
above.  Corrections additionally multiply by a clipped vector estimate
of the prescribed $U$-cutoff $g_0(y)$, queried with a common broad
$Y_f$-domain tuple.  The unnormalized preparations $hR$ and $hRT$
remain on their broad $Y_a$ domain; their usage parameter $t$ specifies
the conditional slice on which they are averaged.  The resulting
second raw signal and the normalized main output carry the $Y_t$ mask.
Dense bins and their preparations and detector tables have
no additional incoming multiplier.
Because $c$ takes values in $[1/2,1]$, every such pre-incoming row
function has a globally Lipschitz $s$th root as a function of its
saturated first entry, first rank, and scaled second entry.  The constants
may depend on the fixed $\zeta$, but not on the stage, rates, or budget.

For clarity, Table~\ref{tab:atom-compression-thresholds} records the
thresholds needed by the compression argument.  For a pre-incoming row
function $H$, its entry factor vanishes through $e_H$, and its row gate
vanishes from rank $g_H$ onward.  Every listed pair satisfies
$3g_H<e_H$, including the broadest witness, for which
$82\vartheta/1000-3(27\vartheta/1000)=\vartheta/1000>0$.
All gates are full in a neighborhood of rank zero.  The correction
preparation in this table is its pre-incoming factor $hR$; its incoming
gate is assembled afterward.

\begin{table}[htbp]
\centering
\begin{tabular}{@{}lll@{}}
\toprule
Pre-incoming row function & $e_H$ & $g_H$ \\
\midrule
Dense bin, including routed bins & $1/2$ & $1/1000$ \\
Dense preparation $hR$ & $1/4$ & $1/1000$ \\
Correction clipped-ratio function & $\theta_1$ & $\theta_1/500$ \\
Correction preparation $hR$ & $\theta_1/2$ & $\theta_1/500$ \\
Detector score & $\delta_0$ & $\delta_0/500$ \\
Auxiliary score $I_{\vartheta}$ & $\vartheta/8$ & $\vartheta/50$ \\
Witness $J_{\vartheta}$ & $82\vartheta/1000$ & $27\vartheta/1000$ \\
\bottomrule
\end{tabular}
\caption{Entry and row-gate thresholds. The entry factor vanishes through
$e_H$, and the row gate vanishes from rank $g_H$ onward. Here
$\vartheta\in\{1/2,\theta_1,\delta_0\}$ as appropriate.}
\label{tab:atom-compression-thresholds}
\end{table}

The bounded first-entry ranges and gate transition intervals already
specified give fixed sensitivity ranges for these functions.  If $w$
is a computed second raw entry, its only use is through $c(w/\zeta)$;
clipping $w$ to $[\zeta/2,\zeta]$ leaves this quantity unchanged.
Its denominator stays in $[1/2,1]$ even for negative or inaccurate
inputs.  Usage masks only remove entries and do not change any of these
support or continuity bounds.

The last access issue is the number of exposed incoming entries. Merely keeping the first $K_c$ would depend on an arbitrary selection and could change nonzero values. Instead, the construction discards every precursor access at an overfull column. The incoming gate has already made all those values zero, so this deletion preserves the table.

\begin{lemma}[Port contract]
\label{lem:atom-ports}
All main and preparatory row tables, auxiliary scores, and the smooth
witnesses just constructed have bounded values and bounded row support.
Every row access class used for a main or preparatory output, or for
access to its auxiliary score, has bounded size and a smooth witness
equal to one on every exposed entry, including exposed entries of weight
zero.  A witness itself is needed as an analytical row function, without
an additional requirement to expose its support through a further
witness.  For a correction or its preparatory table let $A$ be its
pre-incoming class from Equation~\eqref{eq:atom-exposed-class}.  Retain an
exposed entry $(x,y)\in A$ only when
\begin{equation}
 \#\{x':X_u(x')=1,\ (x',y)\in A\}\le K_c.
 \label{eq:atom-port-cap}
\end{equation}
This deletion preserves all table values and gives at most $K-1$
outgoing and at most $K_c$ incoming exposed entries.

These assertions hold separately for every computed version, provided
that its values, ranks, scores, and access classes use the same final
raw arrays and common task tuple.  Access ports denote distinct entries:
the implementation must expose each pair at most once and pair its
incoming and outgoing ports consistently.  Access to auxiliary scores
alone has only the asserted row bound.
\end{lemma}

\begin{proof}
The row and witness assertions are Lemma~\ref{lem:atom-margins}; clipping
and multiplication by bounded factors do not enlarge supports.
If a column has more than $K_c$ entries in $A$, its auxiliary score is
one at all those entries, so its $K_c$th auxiliary rank is one.
Equation~\eqref{eq:atom-incoming-gate} makes $T_y=0$ there.  Consequently
every value deleted by Equation~\eqref{eq:atom-port-cap} was already zero.
The surviving relation has the stated two degree bounds even when zero
weights are retained.  The argument applies to arbitrary raw arrays, so
it applies equally to the estimated and true versions.

The cap removes every precursor port at an overfull column; it does
not choose $K_c$ of them.  It is applied even when the particular entry
has weight zero.  Every surviving port still lies in $A$, so its
pre-incoming witness $I$ equals one regardless of a zero usage mask,
incoming gate, or column multiplier.  The broader auxiliary class
$A_I$ is used to enumerate scores for an incoming rank.  It is not the
capped correction relation used as a paired left factor in a product.

The classes are relations on row and column IDs, rather than multisets
of paths.  Choosing one port per relation element gives the stated
combinatorial interface.  Establishing its single-address enumeration
and port pairing within the space bound is an implementation obligation;
the degree argument does not itself grant random access to both IDs.
\end{proof}

The access requirements can now be summarized without conflating them with accuracy. Every query uses fixed table parameters and one supplied environment. Row ranks include the prescribed dummy coordinates. An exposed port denotes one distinct real entry and lies where its smooth witness is one. Correction and preparatory accesses obey the incoming cap. Values, scores, incoming ranks, and multipliers are all assembled from that same table. In particular, the path used to discover an entry does not change the table at which it is evaluated.  Only completed earlier output types, or stated fixed local
functions and aggregations of those types, may be used as raw inputs of
a later task.  No prefix still waiting for assembly is such an output.
These requirements distinguish finite-support analytical bounds from
the access properties used in the subsequent implementation.
Completing a raw array here means defining its entries from the allowed
earlier outputs, rather than storing the whole array at once.  Its
auxiliary scores and precursor classes are defined before its incoming
gate and cap, so their definitions do not call that same gate or cap.

These contracts apply directly to each computed array after compression
and clamping.  The later comparison in Lemma~\ref{lem:buf-clamping}
preserves relevant row-function values under its stated common-increment
hypothesis.  It need not preserve the finite-accuracy access descriptors
of two calculations.  Each calculation uses its own consistent classes,
and Lemma~\ref{lem:atom-ports} gives their support and degree bounds
separately.

\subsection{How errors propagate through table assembly}

It remains to control how errors in the row functions, incoming scores,
and detector vectors affect an assembled table.  The incoming degree
bound lets us sum these errors by columns without multiplying by the
number of possible source vertices.

We measure error by an eighth moment averaged over environments and source vertices. The source-retention weight compensates for sampling that source. The destination sum stays explicit because bounded support and incoming degree, rather than a second normalization by the universe size, will control it. Root discrepancies are convenient for products of nonnegative gates: the smooth switches were chosen precisely so that their eighth roots vary continuously at zero.

For a real row array $Z$ define
\begin{equation}
 \|Z\|_{u,a;s}^s=
 \frac1N\mathbb E_\sigma\sum_x\frac{X_u(x)}u
                 \sum_{y:Y_a(y)=1}|Z(x,y)|^s,
 \qquad
 \mathfrak d_{u,a}(H,H')=
   \|H^{1/s}-(H')^{1/s}\|_{u,a;s}.
 \label{eq:atom-row-norm}
\end{equation}
Here $\sigma$ has the product law of the rank channels and independent
auxiliary bits specified in Section~\ref{sec:env}.  All table values are
extended by zero outside their membership domain.  The denominator $u$
is the nominal rate, and actual-rate comparability in
Lemma~\ref{lem:env-ranks} absorbs the resulting absolute constants.
For a vertex array on a channel with rate $u$, write
\begin{equation}
 \|v\|_{u;s}^s=
    \frac1N\mathbb E_\sigma\sum_x\frac{X_u(x)}u|v(x)|^s;
 \label{eq:atom-vertex-norm}
\end{equation}
the indicated role channel replaces $X$ when necessary.  These are also
the comparison norms for two sample-dependent versions.  All required
pre-incoming functions and smooth support witnesses are included among
the row outputs whose discrepancies are to be bounded.

Bounded row support and range imply a uniform bound on these row-output
norms, since $|V_l|\le N$ and $\mathbb E X_u/u=O(1)$.  Bounded vertex
arrays have the analogous uniform bound.  Ordinary error is controlled
by root error on bounded nonnegative outputs, since
\begin{equation}
 |z-w|\le s C^{(s-1)/s}|z^{1/s}-w^{1/s}|
 \quad(0\le z,w\le C).
 \label{eq:atom-ordinary-from-root}
\end{equation}

\begin{lemma}[Assembly comparisons]
\label{lem:atom-assembly}
Compare two versions with the same environment, rates, and task tuple.
For a sparse normalized output write $H=P T g$ and
$H'=P'T'g'$, where $P,P'$ are the pre-incoming clipped-ratio row
functions (including their usage masks), $I,I'$ are their auxiliary
scores, and $g,g'$ are their column multipliers.  If $g,g'$ are the
prescribed smooth cutoffs of bounded vertex arrays $U_i,U_i'$, then
\begin{equation}
 \mathfrak d_{u,f}(H,H')^s
 \le C\mathfrak d_{u,f}(P,P')^s
    +C\mathfrak d_{u,f}(I,I')^s
    +C\frac fu\sum_i\|U_i-U_i'\|_{f;s}^s.
 \label{eq:atom-assembly-bound}
\end{equation}
For a sparse preparation omit the last term and use its unnormalized
pre-incoming row function.  The corresponding dense and detector row
functions have no incoming or column-multiplier assembly cost.

If $V_y,V_y'$ are the detector trials formed from
$\mathcal S_i,\mathcal S_i'$ by their column order statistic, then
\begin{equation}
 \|V-V'\|_{f;s}^s
   \le C\mathfrak d_{f,f}(\mathcal S_i,\mathcal S_i')^s.
 \label{eq:atom-detector-comparison}
\end{equation}
\end{lemma}

\begin{proof}
Separate the change in each of the three factors. Then sum only over entries where at least one of the two final outputs is nonzero. This is where the bounded incoming degree can be used. All factors are bounded. On the union $E$ of the two final nonzero
supports, the triangle inequality for the products of $s$th roots gives
\[
 |H^{1/s}-(H')^{1/s}|
 \le C|P^{1/s}-(P')^{1/s}|
     +C|T^{1/s}-(T')^{1/s}|
     +C|g^{1/s}-(g')^{1/s}|.
\]
Outside $E$ the left side is zero.  The pre-incoming discrepancy may be
charged on its entire row domain; it needs no incoming cap.  Each column
has at most $2K_c$ edges in $E$, by Lemma~\ref{lem:atom-ports}.
By Equation~\eqref{eq:atom-rank-lipschitz}, the root-Lipschitz incoming
cutoff, and Equation~\eqref{eq:atom-ordinary-from-root},
\[
 |T_y^{1/s}-(T_y')^{1/s}|^s
 \le C\sup_{x:X_u(x)=1}
            |I(x,y)^{1/s}-I'(x,y)^{1/s}|^s.
\]
Consequently its integrated contribution is at most
\[
 \frac C{Nu}\mathbb E\sum_{y:Y_f(y)=1}
              \sum_{x:X_u(x)=1}
              |I(x,y)^{1/s}-I'(x,y)^{1/s}|^s,
\]
which is $C\mathfrak d_{u,f}(I,I')^s$.

For the multiplier term, the same error at a destination may be encountered from several sources. The incoming cap bounds that multiplicity, and the ratio of the two domain weights is exactly $f/u$. The maximum over a fixed number of $U_i$ values is Lipschitz, and the
multiplier cutoff has Lipschitz root.  Thus the last discrepancy is
bounded in $s$th power by $C\sum_i|U_i(y)-U_i'(y)|^s$.
Summing over $E$ costs at most $2K_c/u$ per retained column.  The vertex
metric has weight $1/f$ instead, giving precisely the factor $f/u$ in
Equation~\eqref{eq:atom-assembly-bound}.  All columns in this comparison
are in the common $Y_f$ domain, including those whose final values carry
a narrower usage mask.  The common-tuple condition ensures that the
vertex discrepancy at $y$ is the same quantity for every old start.
Explicitly, if $m_E(y)$ counts the edges of $E$ entering $y$, then for
each $i$ its contribution obeys
\[
 \frac1{Nu}\mathbb E\sum_{y:Y_f(y)=1}
       m_E(y)|U_i(y)-U_i'(y)|^s
 \le 2K_c\frac fu\|U_i-U_i'\|_{f;s}^s.
\]
This is the column-counting step; it uses no independence between the
two row tables and the vector estimates.

For a detector trial, Equation~\eqref{eq:atom-rank-lipschitz} bounds the
trial discrepancy by the supremum of the input score discrepancies over
retained starts.  Its $s$th power is at most the sum of their $s$th powers.
The factor $Y_f(y)/f$ in the vertex metric is exactly the factor obtained
by reversing the two sums in the row metric with $u=a=f$.
Equation~\eqref{eq:atom-ordinary-from-root} gives
Equation~\eqref{eq:atom-detector-comparison}.
\end{proof}

The required task types are therefore the dense bins $\mathcal B_b$ of
$C_l$, the routed bins $\mathcal B_b^{(j)}$ of $D_l$, detector scores
$\mathcal S_i$, and the pilot and actual correction tables $\mathcal F$,
together with one earlier unnormalized preparation for each normalized
type.  Their true conditional-mean identities are
Lemmas~\ref{lem:atom-dense} and~\ref{lem:atom-correction}.  Their computed
versions use the same bounded row functions and the assembly contract
above.  The next sections bound their discrepancies through the
additive schedule and implement their accesses without storing two
unrestricted vertex addresses.

